% Capitulo: Governança de dados na era da IA
\chapter{Governança de dados na era da IA}
\label{cap:governanca}

\rascunho{
  \textbf{Ideia do capítulo.} Governança de dados sempre teve fama de burocracia:
  comitês, planilhas de dicionário de dados e um data steward fazendo papel de fiscal.
  Com IA, boa parte do trabalho braçal (documentar, classificar, conciliar, avisar) pode
  ser proposto pela máquina, e o data steward passa a ser tutor e revisor do que a IA
  faz. E construir uma aplicação interna que concilia dados e avisa os interessados
  ficou ao alcance de qualquer time com alguém que programe em Python.
  \begin{itemize}
    \item Continuação direta do capítulo~\ref{cap:master-data}; o
      capítulo~\ref{cap:aplicacao-propria} é o guia para construir a própria aplicação.
    \item Retomar a tensão entre governança e liberdade do self-service
      (\veja{sec:governanca}).
    \item Referências de base: \textcite{dama2017}, \textcite{eryurek2021} e, para a ideia
      de governança que não atrapalha, \textcite{seiner2014}.
  \end{itemize}
}

% ---------------------------------------------------------------------------
\section{O que é governança de dados}
\label{sec:gov-o-que-e}

\rascunho{
  \begin{itemize}
    \item O conjunto de papéis, regras e processos que garantem que o dado é confiável,
      bem descrito, protegido e usado de forma adequada \parencite{dama2017}.
    \item Os papéis clássicos: dono do dado (decide as regras do domínio), data steward
      (cuida da qualidade e da documentação no dia a dia), custodiante (cuida da parte
      técnica). \pendente{conferir nomenclatura dos papéis no DMBOK}
    \item O que entra: catálogo e metadados, qualidade, master data, acesso e
      privacidade, linhagem, ciclo de vida.
  \end{itemize}
}

\pendente{escrever}

% ---------------------------------------------------------------------------
\section{Como a governança era vista}
\label{sec:gov-antes}

\rascunho{
  \begin{itemize}
    \item Comitês, políticas longas que ninguém lia, dicionários de dados em planilha
      desatualizados no dia seguinte.
    \item O data steward como fiscal: alguém que mantinha tudo à mão e cobrava dos outros.
      Trabalho pesado, pouco reconhecido, sempre atrasado em relação aos dados.
    \item Por isso a governança costumava perder para o prazo: era a primeira coisa
      cortada quando o projeto apertava.
    \item A proposta de governança não invasiva: formalizar responsabilidades que as
      pessoas já exercem, em vez de criar uma estrutura nova \parencite{seiner2014}.
  \end{itemize}
}

\pendente{escrever}

% ---------------------------------------------------------------------------
\section{O data steward como tutor e revisor da IA}
\label{sec:steward-tutor}

\rascunho{
  \begin{itemize}
    \item O que a IA pode propor: descrições de tabelas e colunas, classificação de dado
      sensível (dado pessoal, financeiro), pares suspeitos de master data, inconsistências
      entre fontes, regras de qualidade a partir dos dados.
    \item O que o steward faz: revisa, aprova, corrige, rejeita. Cada correção vira
      exemplo ou regra, e a próxima proposta melhora (Figura~\ref{fig:governanca-ciclo}).
    \item A mudança de papel: de quem faz tudo à mão para quem ensina e confere. O mesmo
      movimento do TDD (\veja{sec:tdd}) e da documentação como contexto
      (\veja{sec:metadados-contexto}): a IA tornou barato o que antes ficava para depois.
    \item Relato: onde eu já vi ou apliquei isso. \pendente{escrever com relato, sem nomes}
  \end{itemize}
}

\pendente{escrever}

\begin{figure}[htbp]
  \centering
  % Ciclo de governanca com o data steward como tutor e revisor da IA (cap. 14).
\begin{tikzpicture}[diagrama, x=1mm, y=1mm,
    etapa/.style={caixa, text width=34mm, minimum height=13mm, font=\scriptsize}]
  \node[destaquecaixa, etapa] (prop) at (55,28)
    {\textbf{A IA propõe}\\descrições, classificação de dado sensível, pares suspeitos};
  \node[etapa] (rev) at (96,10)
    {\textbf{O steward revisa}\\aprova, corrige ou rejeita};
  \node[etapa] (reg) at (55,-8)
    {\textbf{A decisão é registrada}\\auditoria, correspondência, catálogo};
  \node[etapa] (apr) at (14,10)
    {\textbf{Regras e exemplos melhoram}\\cada correção ensina a próxima proposta};
  \draw[seta] (prop.east) to[bend left=20] (rev.north);
  \draw[seta] (rev.south) to[bend left=20] (reg.east);
  \draw[seta] (reg.west) to[bend left=20] (apr.south);
  \draw[seta] (apr.north) to[bend left=20] (prop.west);
  \node[rotulo] at (55,10) {o steward como\\tutor e revisor};
\end{tikzpicture}

  \caption{O ciclo de governança com a IA propondo e o data steward revisando. Cada
    decisão registrada melhora a próxima proposta.}
  \label{fig:governanca-ciclo}
  \fonte{elaborado pelo autor.}
\end{figure}

% ---------------------------------------------------------------------------
\section{Uma aplicação interna ao alcance de qualquer time}
\label{sec:gov-aplicacao-interna}

\rascunho{
  \begin{itemize}
    \item Exemplo concreto: um programa em Python, de uso interno, que roda depois de cada
      carga, encontra a mesma dimensão com nomes diferentes (o mesmo cliente escrito de
      três jeitos) e manda uma mensagem para o dono do domínio revisar.
    \item Antes isso exigia um projeto; hoje cabe em dias, com a IA ajudando a escrever
      o código. O difícil continua sendo combinar com o negócio quem decide e o que
      conta como \enquote{o mesmo cliente}.
    \item O capítulo~\ref{cap:aplicacao-propria} mostra o que considerar para levar isso
      a sério, com um trecho de código (\veja{sec:motor-conciliacao}).
  \end{itemize}
}

\pendente{escrever}

% ---------------------------------------------------------------------------
\section{Ferramentas de mercado}
\label{sec:gov-ferramentas}

\rascunho{
  \begin{itemize}
    \item A Tabela~\ref{tab:ferramentas-governanca} organiza as opções por categoria. É
      um retrato de 2026: nomes, donos e preços mudam rápido, então o texto deve falar
      das categorias e usar os nomes só como exemplo.
    \item Movimentos recentes que valem uma frase: a Salesforce concluiu a compra da
      Informatica no fim de 2025, e a SAP concluiu a compra da Reltio em 2026.
      \pendente{conferir datas nas fontes oficiais antes da impressão}
    \item Plataformas nativas (catálogo do próprio armazém de dados) resolvem boa parte
      dos casos de quem concentra os dados numa plataforma só; catálogos independentes
      fazem mais sentido com muitas plataformas ou em setores regulados.
    \item Nada impede a empresa de construir a própria aplicação, sobretudo para começar
      pequeno (capítulo~\ref{cap:aplicacao-propria}).
  \end{itemize}
}

\pendente{escrever}

\begin{table}[htbp]
  \centering
  \caption{Ferramentas de governança, catálogo e master data por categoria.}
  \label{tab:ferramentas-governanca}
  \small
  \begin{tabularx}{\textwidth}{@{}>{\raggedright\arraybackslash}p{0.25\textwidth}
      >{\raggedright\arraybackslash}X>{\raggedright\arraybackslash}X@{}}
    \toprule
    Categoria & Exemplos & Quando faz sentido \\
    \midrule
    Catálogo e governança independentes & Alation, Collibra, Atlan, Informatica (Salesforce) & Muitas plataformas de dados, setor regulado, programa corporativo \\
    \addlinespace
    Nativas da plataforma & Databricks Unity Catalog, Snowflake Horizon, Microsoft Purview & Dados concentrados numa plataforma só \\
    \addlinespace
    Código aberto & DataHub, OpenMetadata & Time com engenharia para hospedar e manter \\
    \addlinespace
    Master data & Informatica MDM (Salesforce), Reltio (SAP), Profisee, Semarchy, Stibo Systems & Vários domínios, muitas fontes, regras complexas de conciliação \\
    \addlinespace
    Aplicação própria & Python, um banco de dados e um canal de mensagens & Começar pequeno, um domínio por vez, regras muito específicas do negócio \\
    \bottomrule
  \end{tabularx}
  \fonte{elaborado pelo autor com base nos sites dos fornecedores, em novembro de 2026.
    \pendente{conferir a tabela antes da impressão}}
\end{table}

% ---------------------------------------------------------------------------
\section{O que não se delega à IA}
\label{sec:gov-nao-delegar}

\rascunho{
  \begin{itemize}
    \item Quem é dono de cada domínio e quem responde por ele.
    \item As regras de negócio: o que é um cliente ativo, quando dois cadastros são o
      mesmo.
    \item Quem pode ver o quê, e as obrigações legais com dado pessoal (LGPD).
      \pendente{citar a LGPD}
    \item A decisão final em caso de dúvida: a IA propõe, uma pessoa decide e assina.
  \end{itemize}
}

\pendente{escrever}

% ---------------------------------------------------------------------------
\begin{leituras}
  \leitura{eryurek2021}{Uma visão atual e prática de governança, com atenção a dados na
    nuvem, qualidade, privacidade e ao papel de cada pessoa.}
  \leitura{seiner2014}{A ideia de governança que formaliza o que as pessoas já fazem, em
    vez de criar uma burocracia nova. Combina bem com o steward como revisor.}
  \leitura{dama2017}{O guia de referência para o vocabulário e para as áreas da gestão de
    dados.}
\end{leituras}
